\documentclass[10pt,a4paper]{amsart}
\usepackage[margin=19mm]{geometry}
\usepackage{amsmath,amssymb,mathtools}
\usepackage[scr=rsfs,cal=euler]{mathalfa}
\usepackage{xfrac}
\usepackage[unicode,hidelinks,bookmarks=false]{hyperref}
\newcommand{\ie}{i.e.\ }
\newcommand{\cf}{cf.\ }
\newcommand{\resp}{respectively\ }
\newcommand{\Subsection}{\subsection}
\providecommand{\nobreakdash}{\nobreak\hbox{-}}
\setlength{\emergencystretch}{2em}
\multlinegap0pt
\usepackage{bm}
\usepackage{bbm}
\usepackage{dsfont}
\usepackage{xspace}
\usepackage{cancel}
\usepackage{mathtools}
\usepackage{amsthm}
\makeatletter
\@ifundefined{prop}{\newtheorem{prop}{Proposition}}{}
\makeatother
\newcommand{\GO}{\mathrm{GO}}
\newcommand{\wQ}{\widehat{q}}
\newcommand{\wT}{\widehat{\tau}}
\title[Witt rings, Pfister forms, and equivariant birational geomet]{Witt rings, Pfister forms, and equivariant birational geometry}
\author{Brendan Hassett, Yuri Tschinkel}
\date{Source: arXiv:2608.17821v1 (CC BY 4.0)}
\begin{document}
\maketitle

\noindent\emph{Source and licence.} Witt rings, Pfister forms, and equivariant birational geometry. Version arXiv:2608.17821v1 (CC BY 4.0), distributed under the Creative Commons Attribution 4.0 International licence, as recorded in the arXiv licence metadata for this exact version and shown on the abstract page https://arxiv.org/abs/2608.17821v1. Full rights notice: licenses/arxiv-2608.17821-CC-BY-4.0.txt.\par\smallskip

\noindent\textit{Editorial scope.} Counted unit: the Prop whose source label is \texttt{None} in the article (source0.tex lines 2502--2521, source label \texttt{None}), delivered together with the article's own complete original proof of it (source0.tex lines 2522--2582, one \texttt{proof} environment of 61 lines). No mathematics is written, repaired or re-derived: statement and proof are the article's own text line for line. Required context retained uncounted: none. Every definition, display and label of those lines is retained unchanged. Omitted from this package and not redistributed: 1--2501, 2583--2605. Nothing inside a retained excerpt is omitted or altered: every retained body is delivered exactly as the article's lines are written, so no omission receipt of any kind accompanies this package. Citations: the retained excerpts carry no citation command at all, so no bibliography entry of the article is transcribed and no reference is redistributed. Reference adaptation: none was needed and none was made; every reference inside the delivered unit resolves inside it, so no unresolved reference or citation is produced. Printed slips kept literal: no printed word, symbol, hypothesis or number of the counted statement or of its proof was altered. Layout and class: the article's own class is \texttt{amsart} with packages including amsmath, amssymb, wasysym, graphicx, xcolor, hyperref, xy; the wrapper is the lane's pinned \texttt{amsart} editorial wrapper at 10pt on A4, which restates the theorem-like environments the retained text needs and adds the article's own macro definitions that the retained text needs (\texttt{\textbackslash GO}, \texttt{\textbackslash wQ}, \texttt{\textbackslash wT}), with extra wrapper packages (bm, bbm, dsfont, xspace, cancel, mathtools). The delivered numbering is the unit's own. Assets: no retained span contains a figure environment, a graphics-inclusion command, a PDF-page inclusion, a table or a TikZ picture; no image file is copied into this package, the package declares no asset dependency and no image file is redistributed with it. Licence: version arXiv:2608.17821v1 of this article is distributed under the Creative Commons Attribution 4.0 International licence, as recorded in the arXiv licence metadata for this exact version and shown on the abstract page https://arxiv.org/abs/2608.17821v1. A full rights notice is delivered at licenses/arxiv-2608.17821-CC-BY-4.0.txt. Characters: every retained excerpt is pure ASCII, so the delivered engine, XeLaTeX with its default Latin Modern text font, reports no missing character.
\begin{prop}
Assume $(V,q)$ is non-degenerate, $G$-invariant, and 
equivariantly multiplicative with rational map
$$\tau:V \dashrightarrow \GO(V,q)$$
that is odd as a function of $V$.
Let $W=\operatorname{span}(1,\alpha)$ be the regular representation of $C_2$, with $\alpha$ a non-trivial character of $C_2$. Write
$V\otimes W$ for the resulting representation of $G\times C_2$ and 
$$\wQ(x\otimes 1+x'\otimes \alpha) = q(x) - q(x')$$
for the induced quadratic form.

Then this is equivariant multiplicative for $G\times C_2$, i.e.~there is a linear transformation
$$\wT(x\otimes 1 + x'\otimes \alpha): V\otimes W \rightarrow V\otimes W,$$
with coefficients rational functions in $x$ and $x'$,
with
$$\wQ(\wT(x\otimes 1 + x'\otimes \alpha)\cdot
(y\otimes 1 + y'\otimes \alpha))=
\wQ(x\otimes 1 + x'\otimes \alpha)
\wQ(y\otimes 1 + y'\otimes \alpha).$$
Moreover, $\wT$ is odd as a function on $V\otimes W$.
\end{prop}

\begin{proof}
Set 
\begin{equation}
\wT(x \otimes 1 + x' \otimes \alpha)
= \left(
\begin{matrix} \tau(x) & -\tau(x') \\
        -\tau(x') & \frac{q(x)}{q(x')}\tau(x')\tau(x)^{-1}\tau(x')
        \end{matrix} \right)
\label{eqn:wT}
\end{equation}
a matrix expressed in terms of $\{1,\alpha \}.$
We expand out 
$$\wQ\left(\wT(x\otimes 1 + x'\otimes \alpha)\cdot
(y\otimes 1 + y'\otimes \alpha)\right),$$
using the functional relation along with the fact
that $q$ is quadratic, as follows:
\begin{align*}
q(\tau(x)y - \tau(x')y') - 
q(-\tau(x')y+\frac{q(x)}{q(x')}\tau(x')\tau(x)^{-1}\tau(x')y')  \\
= 
q(\tau(x)y - \tau(x')y') - 
q(x') q(-y+\frac{q(x)}{q(x')}\tau(x)^{-1}\tau(x')y')\\
=
q(\tau(x)y - \tau(x')y') - 
\frac{q(x')}{q(x)} q(-\tau(x)y+\frac{q(x)}{q(x')}\tau(x')y') \\
= q(\tau(x)y - \tau(x')y') - \frac{1}{q(x)q(x')}
q(-q(x')\tau(x)y+q(x)\tau(x')y').
\end{align*}
Writing $b(,)$ for the symmetric bilinear form 
associated with $q$, this becomes
\begin{align*}
&q(\tau(x)y) - 2b(\tau(x)y,\tau(x')y')+q(\tau(x')y') \\
&-\frac{q(q(x')\tau(x)y)-2b(q(x')\tau(x)y,q(x)\tau(x')y') + q(q(x)\tau(x')y')}{q(x)q(x')}.
\end{align*}
The bilinear terms cancel, so clearing denominators
yields the simplification
$$q(\tau(x)y) + q(\tau(x')y') - q(\tau(x)y)\frac{q(x')}{q(x)}
- q(\tau(x')y') \frac{q(x)}{q(x')}.$$
Applying the functional relation again gives
$$q(x)q(y) + q(x')q(y') - q(y)q(x') - q(y')q(x)
=(q(x)-q(x'))(q(y)-q(y')),$$
as desired.

Equivariance of $\wT$ under the action of $G$ follows from the equivariance of $\tau$. Equivariance under
$C_2$ follows because $\tau$ is odd:
$$\left(
\begin{matrix} \tau(x) & -\tau(-x') \\
        -\tau(-x') & \frac{q(x)}{q(x')}\tau(-x')\tau(x)^{-1}\tau(-x')
        \end{matrix} \right)
        =
        \left(
\begin{matrix} \tau(x) & \tau(x') \\
        \tau(x') & \frac{q(x)}{q(x')}\tau(x')\tau(x)^{-1}\tau(x')
        \end{matrix} \right); 
$$
this is the conjugation of our original matrix by
$$\alpha = \left( \begin{matrix} 1 & 0 \\ 0 & -1 \end{matrix}
\right).$$
Formula (\ref{eqn:wT}) also shows that $\wT$ is odd 
as a function of $(x,x')$.
\end{proof}

\end{document}
